\documentclass[11pt,a4paper]{article}

\usepackage[T1]{fontenc}
\usepackage[utf8]{inputenc}
\usepackage{lmodern}
\usepackage[a4paper,margin=26mm]{geometry}
\usepackage{amsmath,amssymb,amsthm,mathtools}
\usepackage{booktabs,array}
\usepackage{microtype}
\usepackage{xcolor}
\usepackage{enumitem}
\usepackage{url}
\usepackage[colorlinks=true,linkcolor=blue!55!black,citecolor=blue!55!black,
            urlcolor=blue!55!black]{hyperref}

\newtheorem{theorem}{Theorem}
\newtheorem{lemma}[theorem]{Lemma}
\newtheorem{proposition}[theorem]{Proposition}
\newtheorem{remark}[theorem]{Remark}
\newcommand{\F}{\mathbb{F}}
\newcommand{\K}{\mathbb{K}}
\newcommand{\E}{\mathbb{E}}
\newcommand{\Prb}{\mathop{\rm Pr}}
\newcommand{\Enc}{\mathsf{Enc}}
\newcommand{\Sign}{\mathsf{Sign}}
\newcommand{\Verify}{\mathsf{Verify}}
\newcommand{\KeyGen}{\mathsf{KeyGen}}
\newcommand{\rank}{\mathop{\rm rank}}
\newcommand{\pk}{\mathsf{pk}}
\newcommand{\sk}{\mathsf{sk}}

\title{\bfseries Repeated Reed--Solomon Openings Break Sigurd\\
\large Few-Signature Witness Recovery and Fresh-Message Forgery\\
Against All Three Submitted Parameter Sets}
\author{Martin Feussner\\
\small Selmer Center, University of Bergen\\
\small \texttt{martin.feussner@uib.no}}
\date{25 September 2026}

\begin{document}
\maketitle

\begin{center}
\fbox{\parbox{0.92\linewidth}{\small\textbf{Provenance and review status.}
The attack, derivation, implementation, experiments, and initial manuscript
were produced independently by OpenAI Codex using Daybreak Blue at max
reasoning effort during an AI-run audit of NGCC signature candidates.  At the
date above, no human cryptanalyst has independently verified the argument,
implementation, measurements, or novelty assessment.  This report is a
preliminary disclosure for human review and reproduction.}}
\end{center}
\vspace{0.7em}

\begin{abstract}
Sigurd is an NGCC submission whose signing witness is a regular-syndrome-
decoding vector.  Its vector-commitment layer is intended to hide that witness
by appending fresh random field elements and applying one Reed--Solomon
encoding.  Both submitted implementations instead pad the input to a multiple
of 32 and encode 32 disjoint short chunks.  Most early chunks contain only the
long-term witness and no fresh padding.  A signature explicitly opens selected
encoded-witness symbols, so openings from different signatures repeat within
those chunks.

We give a passive key-recovery attack.  For each witness-only chunk, the
attacker collects enough distinct public evaluations across signatures and
inverts a small Vandermonde matrix.  This recovers nearly the entire one-hot
witness.  Subtracting its contribution from the public syndrome leaves a thin
binary linear system that recovers the remaining blocks.  The recovered
witness is an equivalent signing key and can run the public prover algorithm
on arbitrary new messages.  The attack does not exploit the placeholder hash,
a side channel, malformed inputs, or a rare secret-key event.

End-to-end experiments used the unmodified submitted reference verifiers.  We
recovered and forged for 20 of 20 Sigurd-128 keys with 4--5 signatures, 10 of
10 Sigurd-256 keys with 5--7 signatures, and 5 of 5 Sigurd-512 keys with 7--8
signatures.  Mean query counts were 4.5, 6.4, and 7.2.  Total single-threaded
CPU time, including generation and verification of every queried signature and
an accepted fresh-message forgery, was 0.114--0.137 seconds, 1.077--1.432
seconds, and 14.198--16.180 seconds, respectively, on an Intel Core i5-6500.
Representative peak resident memory was 14.8, 56.3, and 249.5 MB.  Thus every
submitted security level is practically broken on an ordinary laptop.
\end{abstract}

\section{Introduction}

Sigurd is a post-quantum signature candidate submitted to the
Next-generation Commercial Cryptographic Algorithms Program (NGCC).  It proves
knowledge of a regular-syndrome-decoding (RSD) witness using a
VOLE-in-the-head construction and a Reed--Solomon/Merkle vector commitment
\cite{Sigurd}.  The public key contains a binary matrix seed and a syndrome
$y=He$.  The secret witness $e$ is divided into blocks of six bits, with
exactly one bit set in every block.

The vector-commitment layer is central to the claimed zero knowledge.  The
specification appends fresh random padding to the lifted witness and describes
one global Reed--Solomon encoding.  Random padding is supposed to make the few
opened codeword positions statistically hide the unpadded input.  The
submitted C implementations use a different concrete map: they partition the
padded coefficient vector into 32 chunks and encode every chunk independently.
The fresh padding is at the end, so many chunks consist entirely of the
long-term witness.  Their codewords are identical in every signature.

This turns ordinary signatures into an erasure-decoding oracle.  Each
signature serializes the opened encoded-witness values themselves.  Once an
attacker has seen $k$ distinct positions of a length-$s$ chunk, interpolation
recovers all $k$ input symbols.  Applying this independently to the
witness-only chunks recovers all but 9, 17, and 10 one-hot blocks in the three
parameter sets.  The public equation $He=y$ then determines the remaining 54,
102, and 60 bits by Gaussian elimination.

Our contributions are:

\begin{itemize}[leftmargin=1.5em]
  \item We identify a deterministic privacy failure in the concrete
        Reed--Solomon map used by both submitted implementations.
  \item We give an exact, polynomial-time witness-recovery algorithm from a
        handful of ordinary signatures.  It needs neither chosen messages nor
        access to signing randomness.
  \item We prove that $k$ distinct openings in every witness-only chunk plus a
        full-rank residual public matrix suffice for complete recovery.
  \item We turn the recovered witness into an equivalent signing key and
        produce signatures on fresh messages that the unmodified submitted
        verifiers accept.
  \item We validate the attack across 35 independently generated keys and all
        three submitted security levels, with at most eight signatures and
        less than 17 CPU-seconds in every trial.
\end{itemize}

The result is a full EUF-CMA break of the submitted implementations, rather
than signature malleability or a distinguisher.  The abstract construction
with one correctly implemented global Reed--Solomon code is outside the claim
of this attack; Section~\ref{sec:scope} separates the two precisely.

\section{Relevant parts of Sigurd}
\label{sec:sigurd}

We use the submission's notation where possible.  Only the key relation,
witness lifting, vector commitment, and serialized openings are needed for the
attack.

\subsection{The RSD signing witness}

For every submitted level, the public key defines
\[
             H\in\F_2^{r\times m},\qquad y=He\in\F_2^r.
\]
The secret vector is split into $n=m/6$ blocks
\[
            e=(e_0\|\cdots\|e_{n-1}),
            \qquad e_i\in\F_2^6,
            \qquad \operatorname{wt}(e_i)=1.
\]
If the nonzero position of $e_i$ is $a_i\in\{0,\ldots,5\}$, Sigurd lifts the
block to
\[
                       \widehat e_i=Y^{a_i}.
\]
The implementation stores an element of the ambient extension representation
as $\tau$ coefficients over $\K=\F_{2^{16}}$.  Therefore a lifted secret block
is simply a raw coefficient vector with one in coordinate $a_i$, zeros in the
other first six coordinates, and zeros in coordinates $6,\ldots,\tau-1$.

For each signature the signer appends a fresh augmentation tail of length
$T$ to obtain
\[
 \widehat e'=(\widehat e_0,\ldots,\widehat e_{n-1},
                       a_0,\ldots,a_{T-1})\in\E^{N'},
 \qquad N'=n+T.
\]
The public key is reusable and the prefix of length $n$ is fixed; only the tail
is refreshed.

\subsection{What a signature opens}

The signer encodes $\widehat e'$ and a fresh mask vector, places pairs of
encoded symbols in Merkle trees, and derives $Q$ distinct opening positions
from the Fiat--Shamir transcript.  The serialized signature ends with
\[
        \Enc(\widehat e')_{j_0},\ldots,
        \Enc(\widehat e')_{j_{Q-1}},
\]
together with authentication hashes.  The verifier reconstructs the masking
symbols from the public masked word and these encoded-witness values.  Thus the
witness-code symbols are transmitted in the clear.  Their positions are also
public: anyone can replay the transcript hash and the deterministic
without-replacement sampler.

This disclosure is safe only if each queried code symbol is hidden by the
fresh augmentation.  That property holds for the intended global encoding,
but it fails for the submitted chunked encoder.

\subsection{The concrete 32-way encoder}
\label{sec:encoder}

Let
\[
             k=\left\lceil\frac{N'}{32}\right\rceil
\]
and pad $\widehat e'$ with zeros to length $32k$.  Both the reference and
optimized implementations process 32 chunks
\[
       (\widehat e'_{ck},\ldots,\widehat e'_{ck+k-1}),
       \qquad 0\le c<32,
\]
independently.  For each raw coefficient lane $d\in\{0,\ldots,\tau-1\}$ they
form
\begin{equation}
       F_{c,d}(X)=\sum_{b=0}^{k-1}
                    \widehat e'_{ck+b,d}X^b\in\K[X]
\label{eq:chunkpoly}
\end{equation}
and output evaluations at
\[
                x_j=\alpha^{j+1},\qquad 0\le j<s,
                \qquad s=\rho^{-1}k,
\]
where $\alpha=2$ generates $\K^*$.  Global codeword position $cs+j$ contains
all $\tau$ values $F_{c,d}(x_j)$.  The implementation emits the first $N$
positions of the resulting 32 codewords.

This map is a direct sum of short Reed--Solomon evaluation maps.  It is not the
single degree-$<N'$ polynomial evaluation described by the specification.  In
particular, if
\[
                     c<C:=\left\lfloor\frac{n}{k}\right\rfloor,
\]
then all $k$ inputs to chunk $c$ belong to the fixed witness prefix.  No fresh
augmentation symbol enters Equation~\eqref{eq:chunkpoly}, so the entire chunk
repeats under key reuse.

Table~\ref{tab:params} gives the resulting attack dimensions.  The column
$R=n-Ck$ is the number of secret blocks left outside the completely fixed
chunks.

\begin{table}[ht]
\centering
\caption{Submitted parameters and chunk dimensions used by the attack.}
\label{tab:params}
\small
\begin{tabular}{@{}lrrrrrrrrrr@{}}
\toprule
Instance & $m$ & $n$ & $N'$ & $N$ & $Q$ & $\tau$ & $k$ & $s$ & $C$ & $R$ \\
\midrule
Sigurd-128 & 1302 & 217 & 392  & 1568  & 160 & 10 & 13 & 52  & 16 & 9  \\
Sigurd-256 & 2748 & 458 & 651  & 5208  & 171 & 16 & 21 & 168 & 21 & 17 \\
Sigurd-512 & 5676 & 946 & 1239 & 19824 & 256 & 32 & 39 & 624 & 24 & 10 \\
\bottomrule
\end{tabular}
\end{table}

\section{Key recovery}
\label{sec:attack}

The attack has two algebraic stages.  It first interpolates the repeated
witness-only chunks and then completes the short unknown suffix using the
public syndrome.

\subsection{Interpolating a fixed chunk}

For a fixed chunk $c<C$, collect openings at $k$ distinct local indices
$j_0,\ldots,j_{k-1}$.  For each raw lane $d$, the public values satisfy
\[
 \begin{pmatrix}
 F_{c,d}(x_{j_0})\\[-1mm]\vdots\\F_{c,d}(x_{j_{k-1}})
 \end{pmatrix}
 =
 \begin{pmatrix}
 1&x_{j_0}&\cdots&x_{j_0}^{k-1}\\
 \vdots&\vdots&&\vdots\\
 1&x_{j_{k-1}}&\cdots&x_{j_{k-1}}^{k-1}
 \end{pmatrix}
 \begin{pmatrix}
 \widehat e'_{ck,d}\\[-1mm]\vdots\\\widehat e'_{ck+k-1,d}
 \end{pmatrix}.
\]
The evaluation points are distinct, so this Vandermonde matrix is invertible
over $\K$.  One inversion recovers all $\tau$ lanes and hence all $k$ lifted
witness blocks in the chunk.

\begin{lemma}
For every $c<C$, any $k$ distinct opened positions in chunk $c$ uniquely
determine $\widehat e'_{ck},\ldots,\widehat e'_{ck+k-1}$.
\end{lemma}

\begin{proof}
Equation~\eqref{eq:chunkpoly} has degree less than $k$.  Evaluation at $k$
distinct field points is a bijection from its coefficient vector to its value
vector.  The same Vandermonde matrix applies independently to every raw lane.
\end{proof}

After all $C$ chunks have been decoded, the attacker has $Ck$ of the $n$
secret blocks.  The expected number of distinct local positions seen in one
chunk after $t$ independent signatures is
\begin{equation}
             s\left(1-\left(1-\frac{Q}{N}\right)^t\right),
\label{eq:coverage}
\end{equation}
because a fixed global position is opened with probability $Q/N$ in each
signature.  The attacker simply continues until every fixed chunk has at
least $k$ distinct positions; this is observable and requires no guess.

\subsection{Completing the suffix from the public key}

Let $e_K$ denote the recovered $6Ck$ prefix bits and let $e_U$ contain the
remaining $u=6R$ bits.  Partition the public matrix accordingly:
\[
                       H=(H_K\mid H_U).
\]
The public equation gives
\begin{equation}
                       H_Ue_U=y+H_Ke_K.
\label{eq:residual}
\end{equation}
For the submitted levels, $u$ is only 54, 102, and 60.  Binary Gaussian
elimination therefore recovers $e_U$ immediately.  Each resulting six-bit
block can be checked for weight one, and the complete vector can be verified
against $He=y$.

For a uniformly random $r\times u$ binary matrix with $r\ge u$, the exact
full-column-rank probability is
\[
                    \prod_{i=0}^{u-1}(1-2^{i-r}),
\]
which is overwhelmingly close to one at the submitted dimensions.  If a rare
key gives a rank deficiency, the attacker solves the small affine space and
uses the one-hot conditions; the observed attack is therefore not based on a
special key property.

\subsection{From a witness to a forgery}

The serialized secret key stores a seed that regenerates $e$, so witness
recovery need not reproduce the secret-key byte string.  The signing prover,
however, uses only the expanded public matrix, public syndrome, and witness
vector.  An attacker can lift the recovered $e$, sample fresh augmentation and
masking values, and execute the public signing protocol exactly as an honest
signer.  This is an equivalent signing key.

\begin{theorem}
Suppose the attacker has collected $k$ distinct openings in every chunk
$c<C$, and the residual system in Equation~\eqref{eq:residual} has a unique
solution.  Then the attacker recovers the complete Sigurd signing witness and
can generate an accepted signature on any fresh message in polynomial time.
\end{theorem}

\begin{proof}
The lemma recovers the first $Ck$ lifted blocks exactly.  Their one-hot raw
representations give the corresponding binary prefix $e_K$.  The unique
solution of Equation~\eqref{eq:residual} is the remaining witness suffix
$e_U$, so the concatenation is the complete valid RSD witness.  Sigurd's prover
is a proof of knowledge whose honest algorithm takes that witness and public
instance as input.  Running it with fresh randomness produces a complete
transcript satisfying every verifier equation.  Fiat--Shamir challenge values
are recomputed from the new message and transcript, so the resulting signature
is a fresh-message forgery rather than a replay.
\end{proof}

\section{Complete attack algorithm}

The following procedure uses only the public key and valid signatures under
one key.

\begin{enumerate}[leftmargin=1.8em,label=\textbf{\arabic*.}]
  \item Parse each signature and replay its public Fiat--Shamir hash chain to
        derive the $Q$ distinct opening positions.
  \item Pair each position with the corresponding encoded-witness value in
        the final signature field.  Write the position as $cs+j$, with chunk
        $c$ and local evaluation index $j$.
  \item For every $c<C$, retain one value for each distinct $j$.  Repeated
        observations must agree; disagreement detects a parsing error or a
        signature from another key.
  \item Continue collecting signatures until every fixed chunk contains at
        least $k$ distinct local positions.
  \item For each fixed chunk, choose any $k$ positions, form the Vandermonde
        matrix $(x_j^b)$, and solve it over $\F_{2^{16}}$ for all $\tau$ raw
        lanes.
  \item Decode each recovered lifted value as a one-hot six-bit block.  This
        yields $e_K$.
  \item Expand the public matrix from its public seed, compute the residual in
        Equation~\eqref{eq:residual}, and solve the $r\times 6R$ binary system.
        Check the one-hot conditions and $He=y$.
  \item Use the recovered witness and fresh randomness to run Sigurd's public
        prover on a new message.  Serialize the transcript normally.
\end{enumerate}

No step queries or predicts a random oracle.  The hash is used only to recover
positions that the verifier also derives.  The dominant algebra consists of
small Vandermonde and binary systems.

\section{End-to-end experiments}
\label{sec:experiments}

\subsection{Methodology}

We tested the official NGCC submission archive with SHA-256 digest
\begin{center}
\texttt{443ab14257658d89a60e14ad0d4b7f4df9fcdd06da4f42fcb29bbbfc75a97153}.
\end{center}
The experiments compiled the submitted reference sources with GCC 13.3.0 and
\texttt{-O3} on 64-bit Linux.  The machine had an Intel Core i5-6500 CPU at
3.20 GHz, four physical cores, and 16 GiB of RAM.  Every reported process was
single threaded.

For each independently seeded key, we requested signatures on distinct short
messages and verified every genuine signature before feeding it to the
attacker.  Collection stopped as soon as every witness-only chunk had $k$
distinct symbols.  The attack then interpolated the fixed chunks and solved the
public residual system.  Comparison with the signer-generated witness was
performed only after recovery and was not an attack input.  Finally, a signing
transcript was generated directly from the recovered witness for a message
that had never been queried, and the unmodified submitted verification API was
called on that public key, message, and signature.

The timed region included generation and verification of every queried
signature, parsing, key recovery, construction of the forgery, and verification
of the forgery.  A real passive attacker can omit verification of the genuine
signatures, so these timings are conservative.  Peak memory includes the
submission's expanded public matrices and verifier workspaces.

\subsection{Results}

Table~\ref{tab:results} reports all trials.  Every recovered prefix decoded to
valid one-hot blocks, every residual matrix had full column rank, every
recovered witness matched the test secret exactly, and every fresh-message
forgery was accepted.

\begin{table}[ht]
\centering
\caption{End-to-end key recovery and fresh-message forgery results.}
\label{tab:results}
\small
\begin{tabular}{@{}lrrrrrr@{}}
\toprule
Instance & Keys & Success & Signatures & Mean & CPU seconds & Peak RSS \\
\midrule
Sigurd-128 & 20 & 20/20 & 4--5 & 4.50 & 0.114--0.137 & 14.8 MB \\
Sigurd-256 & 10 & 10/10 & 5--7 & 6.40 & 1.077--1.432 & 56.3 MB \\
Sigurd-512 & 5  & 5/5   & 7--8 & 7.20 & 14.198--16.180 & 249.5 MB \\
\bottomrule
\end{tabular}
\end{table}

The signature counts match the coverage calculation.  For example, after four
Sigurd-128 signatures Equation~\eqref{eq:coverage} predicts 18.3 distinct
symbols per 52-symbol chunk, while only 13 are needed.  For Sigurd-256 after
six signatures it predicts 30.5 of 168 symbols, with 21 needed.  For
Sigurd-512 after seven signatures it predicts 54.3 of 624 symbols, with 39
needed.  The stopping rule depends on the least-covered chunk, which explains
the small variation between keys.

Attack storage for the collected fixed codewords is about 17 KB, 113 KB, and
959 KB for the three levels.  The larger measured resident sets come from
expanding the submitted binary matrix as arrays of machine words and from
running the reference verifier; a streaming public-matrix expansion would
substantially reduce them.

\section{Scope of the security argument}
\label{sec:scope}

The Sigurd specification defines a Reed--Solomon encoder as evaluation of one
polynomial whose coefficient vector contains the full padded input
\cite[Sec.~2.5 and Sec.~4.1]{Sigurd}.  It states that randomized padding hides
the unpadded inputs at queried positions.  In a global encoding, every
evaluation depends on the fresh tail, so openings from separate signatures do
not repeat.

The submitted implementations instead expose a direct sum of 32 encoders.
Random padding can hide only chunks that actually contain padding.  For
$c<C$, the opened value is a deterministic linear function of the secret
prefix alone.  The special-honest-verifier zero-knowledge intuition therefore
fails under key reuse, and the specification explicitly permits at least
$2^{64}$ messages per key while the attack uses at most eight.

This distinction affects the scope of the result:

\begin{itemize}[leftmargin=1.5em]
  \item The reference and optimized implementations for all three submitted
        levels use the vulnerable chunk map, so their valid signatures leak
        the witness and admit end-to-end forgery.
  \item The attack is independent of the concrete placeholder hash and remains
        valid if that hash is replaced, because opening positions and values
        must remain public to the verifier.
  \item The attack does not show that a faithful monolithic $[N,N']$
        Reed--Solomon implementation of the abstract vector commitment leaks
        in the same way.
\end{itemize}

A repair must restore fresh masking to every opened linear form.  Implementing
the specified global encoding is the direct approach.  Alternatively, a
chunked design would need fresh private padding inside every chunk and a new
privacy analysis that covers accumulation across the advertised number of
signatures.  Changing only the query count, Merkle layout, or hash function
does not remove the repeated-codeword leakage.

\section{Independent reproduction procedure}
\label{sec:reproduction}

The following procedure is sufficient to reproduce the result without any
attack artifact from the discovery environment.

\begin{enumerate}[leftmargin=1.8em]
  \item Obtain the official Sigurd submission archive and verify the SHA-256
        digest printed in Section~\ref{sec:experiments}.  Build a reference
        instance and its normal key-generation, signing, and verification
        interfaces.
  \item Generate one key pair.  Obtain valid signatures on successively
        numbered messages under the same key.
  \item Parse the fields in the order specified by the submission:
        Merkle roots; two commitment evaluations; the interactive response
        arrays; six final-response elements; $N'$ masked-word elements; a
        little-endian authentication-hash count and those hashes; and finally
        $Q$ opened encoded-witness elements.
  \item Replay the documented Fiat--Shamir chain.  In particular, derive the
        opening seed from the message hash, the polynomial-response hash, and
        the serialized masked word, then run the submission's public
        without-replacement sampler.  This yields the global position attached
        to every final opened value.
  \item Compute $k=\lceil N'/32\rceil$, $s=\rho^{-1}k$, and
        $C=\lfloor n/k\rfloor$.  For positions below $Cs$, group openings by
        $c=\lfloor j/s\rfloor$ and local index $j\bmod s$.  Stop when each
        group has $k$ distinct indices.
  \item Use evaluation points $x_j=2^{j+1}$ in the submitted
        $\F_{2^{16}}$ representation.  Invert the $k\times k$ Vandermonde
        system for every group and apply the inverse to all raw lanes.  Decode
        the resulting monomials $1,Y,\ldots,Y^5$ as one-hot blocks.
  \item Expand $H$ and $y$ from the public key.  Remove the contribution of
        the recovered prefix, solve Equation~\eqref{eq:residual} over $\F_2$,
        and verify regularity and the complete public syndrome.
  \item Run the normal prover with this recovered witness and fresh randomness
        on a new message.  Acceptance by the unmodified verification API is
        the end-to-end forgery test.
\end{enumerate}

Useful diagnostic checks are deterministic.  A repeated global position in a
fixed chunk must carry the same raw value across signatures.  Interpolated
secret blocks must be one-hot in the first six lanes.  The completed witness
must satisfy $He=y$.  Failure of any check indicates a transcript-layout,
endianness, or position-sampler mismatch rather than a statistical failure of
the attack.

\section{Conclusion}

The submitted Sigurd encoders partition one padded vector into 32 independent
Reed--Solomon codewords.  Fresh padding occupies only late chunks, while the
signature publicly opens symbols from every chunk.  Key reuse therefore
reveals repeated evaluations of secret-only polynomials.  Small Vandermonde
systems recover nearly the whole witness, and the public syndrome supplies the
short suffix.

The result is practical at every submitted level: 4--8 ordinary signatures,
less than 17 seconds of single-threaded CPU time, and less than 250 MB of
measured memory.  The recovered witness signs arbitrary new messages, and the
submitted verifiers accepted all 35 experimental forgeries.  Independent human
review is still required, but the present evidence is a complete,
reproducible key-recovery and forgery chain.

\begin{thebibliography}{9}

\bibitem{Sigurd}
Y.~Deng, Y.~Wei, X.~Zhang, Z.~Wu, Z.~Zhang, and L.~Qin,
``Sigurd: Algorithm Specifications,'' Version 1.0, NGCC submission,
29 June 2026.\\
\href{https://www.niccs.org.cn/niccs/Proposal/Public-Key\%20Cryptographic\%20Algorithms/Round\%201\%20candidates/Sigurd.zip}{Official NGCC submission archive.}

\bibitem{RS}
I.~S. Reed and G.~Solomon,
``Polynomial Codes over Certain Finite Fields,''
\emph{Journal of the Society for Industrial and Applied Mathematics},
vol.~8, no.~2, pp.~300--304, 1960.

\end{thebibliography}

\end{document}
